% Partie : sets-functions-relations
% Chapitre : functions
% Section : kinds

\documentclass[../../../include/open-logic-section]{subfiles}

\begin{document}

\olfileid[fr]{sfr}{fun}{kin}
\olsection{Types de fonctions}

\begin{explain}
Il sera utile de classer certains types de fonctions que nous rencontrerons
fréquemment.

Commençons par les fonctions pour lesquelles chaque élément de l’ensemble
d’arrivée est une valeur prise par la fonction. Elles sont dites
!!{surjective}s et peuvent être représentées comme dans la
\olref{fig:surjective}.

\begin{figure}
  \olasset{assets/diagrams/surjective.tikz}
  \caption{Une fonction !!{surjective} prend pour valeur chaque
    !!{element} de son ensemble d’arrivée.}
  \ollabel{fig:surjective}
\end{figure}
\end{explain}

\begin{defn}[Fonction !!{surjective}]
Une fonction $f \colon A \rightarrow B$ est \emph{!!{surjective}} si et
seulement si $B$ est aussi l’image de~$f$ : pour tout $y \in B$, il existe
au moins un $x \in A$ tel que~$f(x) = y$. En symboles :
\[
  (\forall y \in B)(\exists x \in A)f(x) = y.
\]
On appelle une telle fonction !!a{surjection} de $A$ dans $B$.
\end{defn}

\begin{explain}
Pour montrer que $f$ est !!a{surjection}, il faut montrer que chaque
objet de son ensemble d’arrivée est égal à $f(x)$ pour une certaine
entrée $x$.

Toute fonction \emph{induit} une surjection. Étant donnée
$f \colon A \to B$, définissons $f' \colon A \to \ran{f}$ par
$f'(x) = f(x)$. Puisque $\ran{f}$ est \emph{par définition}
$\Setabs{f(x) \in B}{x \in A}$, la fonction $f'$ est nécessairement
!!a{surjection}.
\end{explain}

\begin{explain}
Toute fonction associe à chaque entrée possible une unique sortie.
Certaines fonctions ont en outre la propriété de ne jamais associer
la même sortie à deux entrées différentes. Elles sont dites
!!{injective}s et peuvent être représentées comme dans la
\olref{fig:injective}.
\begin{figure}
  \olasset{assets/diagrams/injective.tikz}
  \caption{Une fonction !!{injective} n’associe jamais la même valeur
    à deux arguments différents.}
  \ollabel{fig:injective}
\end{figure}
\end{explain}

\begin{defn}[Fonction !!{injective}]
Une fonction $f \colon A \rightarrow B$ est \emph{!!{injective}} si et
seulement si, pour chaque $y \in B$, il existe au plus un $x \in A$
tel que~$f(x) = y$. On appelle une telle fonction !!a{injection}
de $A$ dans~$B$.
\end{defn}

\begin{explain}
Pour montrer que $f$ est !!a{injection}, il faut montrer que,
pour tous les !!{element}s $x$ et $y$ de son domaine,
si $f(x)=f(y)$, alors $x=y$.
\end{explain}

\begin{ex}
La fonction constante $f\colon \Nat \to \Nat$ définie par $f(x) = 1$
n’est ni !!{injective} ni !!{surjective}.

La fonction identité $f\colon \Nat \to \Nat$ définie par $f(x) = x$
est à la fois !!{injective} et !!{surjective}.

La fonction successeur $f \colon \Nat \to \Nat$ définie par $f(x) = x+1$
est !!{injective}, mais n’est pas !!{surjective}.

La fonction $f \colon \Nat \to \Nat$ définie par
\[
  f(x) =
  \begin{cases}
    \frac{x}{2} & \text{si $x$ est pair} \\
    \frac{x+1}{2} & \text{si $x$ est impair.}
  \end{cases}
\]
est !!{surjective}, mais n’est pas !!{injective}.
\end{ex}

\begin{explain}
Nous aurons souvent besoin de fonctions à la fois !!{injective}s et
!!{surjective}s. On les dit !!{bijective}s ; elles se représentent
comme dans la \olref{fig:bijective}. Les !!{bijection}s sont aussi
parfois appelées \emph{correspondances biunivoques}, car elles font
correspondre de manière unique les éléments de l’ensemble d’arrivée
à ceux du domaine.
\begin{figure}
  \olasset{assets/diagrams/bijective.tikz}
  \caption{Une fonction !!{bijective} fait correspondre de manière unique
    les éléments de l’ensemble d’arrivée à ceux du domaine.}
  \ollabel{fig:bijective}
\end{figure}
\end{explain}

\begin{defn}[!!^{bijection}]
Une fonction $f \colon A \to B$ est \emph{!!{bijective}} si et seulement
si elle est à la fois !!{surjective} et !!{injective}.
On l’appelle !!a{bijection} de $A$ dans~$B$ (ou entre $A$ et~$B$).
\end{defn}

\end{document}
